\documentclass[11pt, sans, a4paper]{moderncv}
\usepackage[utf8]{inputenc}
\usepackage[T1]{fontenc}
\usepackage{lmodern}
\usepackage{float}
\usepackage{ragged2e}
\usepackage{amsmath,amssymb}
\usepackage{booktabs}
\usepackage{array}
\usepackage{enumitem}
\usepackage{microtype}
\usepackage[most]{tcolorbox}
\usepackage{hyperref}
\usepackage{catchfilebetweentags}
\usepackage{xr}
\externaldocument{manuscript_marked}

% Citation and reference bridge from manuscript_marked.aux
\makeatletter
\def\extractfirst#1#2#3#4{#1}
\def\bibcite#1#2{\@namedef{b@#1}{\extractfirst#2}}
\let\orig@newlabel\newlabel
\def\newlabel#1#2{}
\InputIfFileExists{manuscript_marked.aux}{}{}
\let\newlabel\orig@newlabel

% Fallback definitions for Overleaf (when manuscript_marked.aux is not generated in the same build)
\providecommand{\fallbackref}[2]{\@ifundefined{r@#1}{\@namedef{r@#1}{{#2}{}}}{}}
\providecommand{\fallbackcite}[2]{\@ifundefined{b@#1}{\@namedef{b@#1}{#2}}{}}

\fallbackref{eq:TM_binding}{1}
\fallbackref{eq:relative_stability}{2}
\fallbackref{fig:Fig1}{1}
\fallbackref{fig:Fig2}{2}
\fallbackref{fig:Fig3}{3}
\fallbackref{fig:Fig4}{4}
\fallbackref{fig:Fig5}{5}
\fallbackref{fig:Fig6}{6}
\fallbackref{fig:Fig7}{7}
\fallbackref{fig:Fig8}{8}

\fallbackcite{ASS2026_166886}{66}
\fallbackcite{Abdel-Hafiez2025}{24}
\fallbackcite{ApplSurfSci2025_161321}{59}
\fallbackcite{Gao2018}{36}
\fallbackcite{Ipaves2025}{7}
\fallbackcite{Liao2022}{6}
\fallbackcite{Luo2020}{35}
\fallbackcite{Mastrippolito2021}{25}
\fallbackcite{Moraes2026}{42}
\fallbackcite{NORSKOV2004}{43}
\fallbackcite{Norskov_2005}{41}
\fallbackcite{PCCP2018_15863}{60}
\fallbackcite{Vacuum2020_109438}{62}
\fallbackcite{Vacuum2020_109742}{61}
\fallbackcite{Webster2018}{21}
\fallbackcite{Zala2024}{40}
\makeatother

\moderncvstyle{classic}
\moderncvcolor{blue}

% Define custom colors
\definecolor{acsblue}{RGB}{0, 65, 130}
\definecolor{revbar}{RGB}{31, 110, 165}
\definecolor{revbg}{RGB}{246, 248, 250}
\definecolor{actionblue}{RGB}{0, 90, 170}
\definecolor{addedgreen}{RGB}{0, 100, 60}

% Custom tcolorbox for reviewer comments
\newtcolorbox{reviewercomment}[1][]{%
  enhanced,
  breakable,
  colback=revbg,
  colframe=revbar,
  arc=1.5mm,
  boxrule=0pt,
  leftrule=3.5pt,
  top=7pt,
  bottom=7pt,
  left=10pt,
  right=10pt,
  fontupper=\small\itshape,
  title={\textbf{\upshape #1}},
  fonttitle=\bfseries\small\color{revbar},
  attach boxed title to top left={yshift=-2.5mm, xshift=3mm},
  boxed title style={boxrule=0pt, colframe=white, colback=white, arc=0mm, top=0pt, bottom=0pt, left=2pt, right=2pt}
}

% Custom tcolorbox for quoted manuscript additions
\newtcolorbox{addedtext}[1][]{%
  enhanced,
  breakable,
  colback=blue!3,
  colframe=acsblue,
  arc=1mm,
  boxrule=0pt,
  leftrule=2.5pt,
  top=5pt,
  bottom=5pt,
  left=8pt,
  right=8pt,
  fontupper=\small,
  title={\textbf{\upshape #1}},
  fonttitle=\bfseries\footnotesize\color{acsblue},
  attach boxed title to top left={yshift=-2.5mm, xshift=3mm},
  boxed title style={boxrule=0pt, colframe=white, colback=white, arc=0mm, top=0pt, bottom=0pt, left=2pt, right=2pt}
}

% Command for formatting author response header
\newcommand{\response}{\vspace{0.4em}\noindent\textbf{\color{actionblue}Response:}\hspace{0.5em}}
\newcommand{\manuscriptchange}[1]{\vspace{0.4em}\noindent\textbf{\color{acsblue}[Changes in Revised Manuscript]:}\par\nopagebreak\vspace{0.15em}#1\vspace{0.4em}}

\name{Prof. Pedro A. S. Autreto}{et al.}
\title{Response to Reviewers' Comments}
\address{Center of Natural and Human Sciences, Federal University of ABC (UFABC)}{Santo André, São Paulo, Brazil}{}
\phone[office]{+55 11 4996-7991}
\email{pedro.autreto@ufabc.edu.br}

\begin{document}

%----------------------------------------------------------------------------------------
%   COVER LETTER TO THE EDITOR
%----------------------------------------------------------------------------------------
\recipient{Prof. Ayan Datta}{Associate Editor\\ACS Applied Energy Materials}
\date{\today}
\opening{Dear Prof. Datta,}
\closing{Sincerely yours,\\On behalf of all authors,}

\makelettertitle

\justify

We would like to thank you and the reviewers for the thorough, constructive, and insightful evaluation of our manuscript entitled \textbf{``Coordination-Dependent Hydrogen Adsorption on Co-, Fe-, and Ni-Functionalized Monolayer $\text{CrCl}_3$''} (Manuscript ID: \textbf{ae-2026-02583t}).

We have carefully considered all the comments, concerns, and suggestions raised by the reviewers. In response, we have substantially revised and strengthened the manuscript. The key revisions and enhancements in the resubmitted version include:

\begin{itemize}[leftmargin=1.5em, itemsep=0.3em]
  \item \textbf{Thermodynamic and Dynamical Stability:} Added explicit physical equations for binding and embedding energies, consolidated stability metrics, and expanded the discussion of $300\text{ K}$ AIMD trajectories alongside benchmark stability references.
  \item \textbf{Electronic Structure, Descriptors, and Magnetic Physics:} Provided in-depth physical rationalization of the local magnetic moments, crystal-field splitting at threefold vs. sixfold coordination, and exchange coupling with the ferromagnetic $\text{CrCl}_3$ host. Clarified the descriptor performance of the occupied $d$-band center ($\varepsilon_d^{\mathrm{occ}}$), calculated theoretical overpotentials ($\eta = |\Delta G_{\mathrm{ads}}|/e$), and positioned the catalysts relative to the exchange current density volcano.
  \item \textbf{Methodological Clarifications:} Elaborated on the justification for spin-polarized PBE, the scope of van der Waals (dispersion) interactions, spin-orbit coupling (SOC) approximations in light-ligand halides, single vs. multi-coverage constraints, and the relation between thermodynamic descriptors and full Volmer/Tafel kinetics.
  \item \textbf{Resolution of Discrepancies:} Explicitly resolved the apparent difference between our calculated pristine chemisorption free energy ($\Delta G_{\mathrm{ads}} = 1.82\text{ eV}$ at the top-Cr site) and weakly interacting surface values reported in earlier literature ($\Delta G_{\mathrm{ads}} = 0.13\text{ eV}$), showing they represent chemically and geometrically distinct adsorption regimes.
  \item \textbf{Presentation and Reference Overhaul:} Unified all notations across the text and figures, clarified Bader charge transfer sign conventions, provided consolidated summary tables of all calculated physical parameters, and systematically corrected all bibliographic entries to strictly comply with the ACS style guide.
\end{itemize}

Below, please find our point-by-point responses to all comments raised by the reviewers. For convenience, all modifications in the marked revised manuscript (\texttt{manuscript\_marked.pdf}) are highlighted in \textbf{\color{blue}blue}.

\vspace{1.5em}
\noindent\rule{\textwidth}{0.6pt}

\newpage

%========================================================================================
%   RESPONSE TO REVIEWER 1
%========================================================================================
\section{Response to Reviewer \#1}

\noindent\textbf{Recommendation:} Major revisions needed as noted.

\vspace{0.5em}
\noindent\textbf{General Comments:}\\
\textit{In this manuscript, the authors have used first-principles method to study the influence of transitions metals on the hydrogen adsorption of $\text{CrCl}_3$. Three transition metals are considered. The authors think that these doped metals can reduce the hydrogen adsorption energy of $\text{CrCl}_3$. Although these results are very interesting, it must be again comment after revision.}

\vspace{0.5em}
\noindent We sincerely thank Reviewer 1 for the careful reading of the manuscript, the positive assessment of our findings, and the constructive recommendations. Below, we provide our detailed point-by-point responses.

\vspace{1em}

%--- R1C1 ---
\begin{reviewercomment}[Reviewer 1 --- Comment 1]
  1). The abstract should highlight the key research points of this paper and should be concisely formulated and well summarized.
\end{reviewercomment}

\response We thank the Reviewer for this helpful suggestion. We agree that the Abstract should foreground the key contributions more effectively. We have rewritten the Abstract to immediately and concisely highlight:
\begin{enumerate}[leftmargin=1.5em, itemsep=0.2em]
  \item The coordination-controlled trade-off between transition-metal substrate binding and hydrogen adsorption thermodynamics on monolayer $\text{CrCl}_3$;
  \item Identification of threefold surface-adsorbed Co as the optimal active motif, delivering near-thermoneutral adsorption ($\Delta G_{\mathrm{ads}} = 0.20\text{ eV}$);
  \item The fundamental mechanistic origin of this performance, revealed through crystal-field splitting, orbital exposure, and distortion--interaction energy partitioning.
\end{enumerate}

\manuscriptchange{\textbf{[Location: \textit{manuscript\_marked.pdf: Pages 1--2, Lines 2--21} \textbar\ \textit{manuscript\_clean.tex: Line 119} \textbar\ \textit{manuscript\_marked.tex: Line 138}]}\\ The revised Abstract has been rewritten and is highlighted in blue in the marked manuscript. The new Abstract reads (excerpt): ``\dots{\itshape \ExecuteMetaData[manuscript_marked.tex]{AbstractTradeoff}}''.}

\vspace{0.8em}

%--- R1C2 ---
\begin{reviewercomment}[Reviewer 1 --- Comment 2]
  2). What the evidence for the choose of these doped elements (Co, Fe and Ni) rather than the other transition metals?
\end{reviewercomment}

\response We appreciate the Reviewer's question. The choice of Co, Fe, and Ni was guided by three distinct physical and chemical criteria, which are now explicitly detailed in the Introduction:
\begin{itemize}[leftmargin=1.5em, itemsep=0.2em]
  \item \textbf{Electronic Structure and Volcano Position:} Co ($d^7$), Fe ($d^6$), and Ni ($d^8$) are late $3d$ transition metals possessing partially filled $d$-shells that hybridize effectively with the $\text{CrCl}_3$ host. This positions them within the optimal, near-thermoneutral activity window on the Sabatier volcano curve where $\Delta G_{\mathrm{ads}}$ can be tuned toward zero by controlling the local coordination environment (N{\o}rskov et al., \textit{J. Electrochem. Soc.} 2005, 152, J23; Moraes et al., \textit{Int. J. Hydrogen Energy} 2026, 212, 153692). In contrast, early $3d$ metals (e.g., Ti, V) bind H too strongly, while late metals with filled $d$-shells (e.g., Cu, Zn) offer limited hybridization.
  \item \textbf{Magnetic Compatibility:} Co, Fe, and Ni carry significant magnetic moments ($1\text{--}4\,\mu_{\mathrm{B}}$), compatible with the $\text{Cr}^{3+}$ ($d^3$, $S=3/2$) Heisenberg ferromagnet ground state, enabling spin--coordination--catalysis coupling studies.
  \item \textbf{Experimental Realizability:} Co, Fe, and Ni are the most commonly reported TM functionalization agents for $\text{CrCl}_3$ and related chromium trihalides in both theoretical studies (Luo et al., \textit{Solid State Commun.} 2020, 321, 114048; Gao et al., \textit{Phys. Status Solidi RRL} 2018, 12, 1800105) and experimental investigations (Mastrippolito et al., \textit{Nanoscale Adv.} 2021, 3, 4756--4766; Abdel-Hafiez et al., \textit{Phys. Rev. B} 2025, 112, 115145).
\end{itemize}

\manuscriptchange{\textbf{[Location: \textit{manuscript\_marked.pdf: Page 4, Lines 72--83} \textbar\ \textit{manuscript\_clean.tex: Line 138} \textbar\ \textit{manuscript\_marked.tex: Line 157}]}\\ An explicit paragraph detailing the rationale for selecting Co, Fe, and Ni has been added to Section~1 (Introduction, highlighted in blue). The added text reads: ``{\itshape \ExecuteMetaData[manuscript_marked.tex]{R1C2Rationale}}''}

\vspace{0.8em}

%--- R1C3 ---
\begin{reviewercomment}[Reviewer 1 --- Comment 3]
  3). If the author considers the effect of metal elements on hydrogen adsorption in $\text{CrCl}_3$, they should not only consider the doping model of alloy elements in $\text{CrCl}_3$, but also the structural model of alloy elements in hydrogen adsorption on $\text{CrCl}_3$.
\end{reviewercomment}

\response We thank the Reviewer for this remark. Indeed, establishing the structural model of the functionalized site before and during H adsorption is the central pillar of our study. For every metal (Co, Fe, Ni), we systematically examined two distinct structural coordination configurations:
\begin{enumerate}[leftmargin=1.5em, itemsep=0.2em]
  \item \textbf{Surface-Adsorbed Motif (Threefold Hollow):} The TM sits above the basal Cl plane, coordinated to three surface Cl atoms, retaining an open, geometrically accessible site for direct H binding.
  \item \textbf{Embedded Motif (Sixfold Intercalated):} The TM penetrates the surface Cl layer, achieving distorted octahedral coordination with six Cl atoms from both the top and bottom Cl sublayers, which chemically shields the TM center.
\end{enumerate}
For all six configurations, full atomic and electronic relaxations were performed before and after H adsorption. The results (Section 3.4, Figure 7) demonstrate that the structural coordination motif is the primary factor dictating $\Delta G_{\mathrm{ads}}$.

\manuscriptchange{\textbf{[Location: \textit{manuscript\_marked.pdf: Page 5, Lines 92--108 \& Page 16, Lines 325--330} \textbar\ \textit{manuscript\_clean.tex: Lines 140, 359} \textbar\ \textit{manuscript\_marked.tex: Lines 159, 378}]}\\ We have clarified the dual-configuration structural modeling strategy in both the Introduction and Section 3.4.}

\vspace{0.8em}

%--- R1C4 ---
\begin{reviewercomment}[Reviewer 1 --- Comment 4]
  4). For structural stability of metal-doped $\text{CrCl}_3$, the authors should be affirmed the structural stability such as thermodynamic and dynamical stabilities. If these compounds are dynamically unstable, there is no mean to study the physical properties. In addition, the authors should be given the physical equation of thermodynamic and dynamically. For the discussion of thermodynamic and dynamical stability, the authors should be cited these references: Applied Surface Science 2025;680: 161321. Physical Chemistry Chemical Physics 2018;20: 15863-15870. Vacuum 2020;181: 109742. Vacuum 2020;179: 109438.
\end{reviewercomment}

\response We thank the Reviewer for emphasizing the importance of stability verification. We fully agree and have incorporated explicit physical definitions, equations, and literature citations:

\noindent\textbf{1. Thermodynamic Stability (Equations 1--2):}
The binding energy of a single TM atom on the $\text{CrCl}_3$ host is defined as:
\begin{equation}
  E_{\mathrm{bind}}^{\mathrm{TM}} = E(\mathrm{TM}/\mathrm{CrCl}_3) - E(\mathrm{CrCl}_3) - E(\mathrm{TM}_{\mathrm{iso}})
\end{equation}
The surface binding energies are strongly negative: $\text{Co} = -2.82\text{ eV}$, $\text{Fe} = -2.60\text{ eV}$, and $\text{Ni} = -3.38\text{ eV}$, confirming spontaneous chemisorption.

The relative thermodynamic stability between embedded and surface-adsorbed motifs is:
\begin{equation}
  \Delta E_{\mathrm{emb-ads}} = E_{\mathrm{emb}} - E_{\mathrm{ads}}
\end{equation}
yielding $-0.73\text{ eV}$ (Co), $-0.76\text{ eV}$ (Fe), and $-1.57\text{ eV}$ (Ni). Thermodynamically at 0~K, the sixfold embedded motif is energetically favored for all three transition metals ($\Delta E_{\mathrm{emb-ads}} < 0$), with Ni exhibiting the largest thermodynamic driving force toward incorporation. Furthermore, embedded Ni ($|E_{\mathrm{bind}}| = 4.95\text{ eV}$) exceeds the cohesive energy of bulk Ni (4.44~eV/atom), confirming full thermodynamic stability against bulk-metal aggregation.

\noindent\textbf{2. Dynamical Stability and Coordination Dynamics (AIMD Simulations):}
Dynamical stability was confirmed via \textit{ab initio} molecular dynamics (AIMD) simulations in the $NVT$ ensemble at $T = 300\text{ K}$ using a Nosé--Hoover thermostat over $5\text{ ps}$ (Figure~3(g--i)):
\begin{itemize}[leftmargin=1.5em, itemsep=0.2em]
  \item \textbf{Monolayer Scaffold Integrity:} The $\text{CrCl}_3$ host lattice remained fully intact throughout all trajectories without bond breaking, severe structural distortion, or thermal desorption.
  \item \textbf{Dynamic Verification of Embedded Stability (Ni):} Consistent with its strong thermodynamic driving force ($\Delta E_{\mathrm{emb-ads}} = -1.57\text{ eV}$), the surface-adsorbed Ni atom undergoes spontaneous migration into the interior of the monolayer during the 300~K trajectory (Figure~3(i)), directly verifying that the sixfold embedded configuration is dynamically accessible and stable at room temperature.
  \item \textbf{Kinetic Trapping of Active Surface Sites (Co and Fe):} For Co and Fe, the adatoms remain predominantly exposed above the basal plane throughout the 5~ps dynamics (Figure~3(g,h)). Although embedding is thermodynamically downhill, the energy barrier required to migrate through the constrained CrCl$_3$ hollow region creates a kinetic trap that preserves the catalytically active, surface-exposed sites at room temperature.
\end{itemize}

\manuscriptchange{\textbf{[Location: \textit{manuscript\_marked.pdf: Page 11, Lines 242--250 \& Pages 20--21, Lines 402--412} \textbar\ \textit{manuscript\_clean.tex: Lines 323, 393} \textbar\ \textit{manuscript\_marked.tex: Lines 342, 412}]}\\ Sections~2 and~3.2 have been updated (highlighted in blue in the marked manuscript) to detail both thermodynamic equations (Eqs.~1--2) and dynamical stability. Section~3.2 explicitly highlights: ``{\itshape \ExecuteMetaData[manuscript_marked.tex]{StabilityHierarchy}}''}

\vspace{0.8em}

%--- R1C5 ---
\begin{reviewercomment}[Reviewer 1 --- Comment 5]
  5). In the structural model, which surface of $\text{CrCl}_3$ did the author select for hydrogen adsorption? What was the basis for this selection? If the adsorption free energy of hydrogen on $\text{CrCl}_3$ is positive, can it still serve as a hydrogen storage material?
\end{reviewercomment}

\response We thank the Reviewer for this clarifying question.

We selected the (001) basal plane of monolayer $\text{CrCl}_3$, i.e., the Cl-terminated surface exposed upon mechanical exfoliation or vapor-phase growth. For a free-standing 2D monolayer, this is the only accessible face, consistent with the $\text{CrCl}_3$ literature (Webster2018, Mastrippolito2021, Zala2024).
Also the positive $\Delta G_{\mathrm{ads}}$ (endergonic relative to $\frac{1}{2}\text{H}_2$) indicates unfavorable H adsorption. An ideal HER electrocatalyst operates near thermoneutrality ($\Delta G_{\mathrm{ads}} \approx 0\text{ eV}$). This work targets HER electrocatalysis exclusively, not hydrogen storage.

\manuscriptchange{\textbf{[Location: \textit{manuscript\_marked.pdf: Page 5, Lines 110--115} \textbar\ \textit{manuscript\_clean.tex: Line 144} \textbar\ \textit{manuscript\_marked.tex: Line 163}]}\\
A new paragraph has been added at the beginning of Section~2 (highlighted in blue): ``{\itshape \ExecuteMetaData[manuscript_marked.tex]{BasalPlaneSelection}}''}

\vspace{0.8em}

%--- R1C6 ---
\begin{reviewercomment}[Reviewer 1 --- Comment 6]
  6). In this article, is the author considering $\text{CrCl}_3$ as a hydrogen storage material or as a hydrogen evolution catalyst? If it is considered as a hydrogen storage material, what is its upper limit for hydrogen storage, and does it have potential for commercial value?
\end{reviewercomment}

\response We thank the Reviewer for this important distinction. We explicitly consider transition-metal-functionalized $\text{CrCl}_3$ as a model electrocatalyst for the \textbf{hydrogen evolution reaction (HER)}, evaluated within the Computational Hydrogen Electrode (CHE) framework. It is not evaluated as a bulk hydrogen storage medium. The manuscript has been revised to eliminate any potential ambiguity regarding storage applications.

\manuscriptchange{\textbf{[Location: \textit{manuscript\_marked.pdf: Page 5, Lines 113--115} \textbar\ \textit{manuscript\_clean.tex: Line 144} \textbar\ \textit{manuscript\_marked.tex: Line 163}]}\\ The clarification has been added to the opening paragraph of Section~2 (Computational Methods), as quoted in Response~5 above.}

\vspace{0.8em}

%--- R1C7 ---
\begin{reviewercomment}[Reviewer 1 --- Comment 7]
  7). In this paper, the author considers the hydrogen adsorption capability of $\text{CrCl}_3$. However, greater attention should be given to its catalytic performance in hydrogen adsorption, such as discussions on the hydrogen adsorption Gibbs free energy, overpotential volcano plots, and exchange current density volcano. For the discussion of hydrogen adsorption Gibbs free energy, overpotential volcano plots, and exchange current density volcano the authors should be cited these references: International Journal Of Hydrogen Energy 2026;265: 156967. Applied Surface Science 2026;737: 166886.
\end{reviewercomment}

\response We thank the Reviewer for this constructive suggestion. In the revised manuscript, we have expanded the catalytic analysis within the CHE framework, $\eta = |\Delta G_{\mathrm{ads}}|/e$. Where the surface-adsorbed motifs: $\eta = 0.20\text{ V}$ (Co), $0.51\text{ V}$ (Fe), $0.88\text{ V}$ (Ni). While the embedded motifs: $\eta > 1.5\text{ V}$. Following the Sabatier principle and Butler--Volmer modeling, $i_0$ peaks at $\Delta G_{\mathrm{ads}} \approx 0\text{ eV}$. Surface-adsorbed Co ($\Delta G_{\mathrm{ads}} = 0.20\text{ eV}$) resides nearest to the volcano apex.

\manuscriptchange{\textbf{[Location: \textit{manuscript\_marked.pdf: Page 28, Lines 527--536} \textbar\ \textit{manuscript\_clean.tex: Line 452} \textbar\ \textit{manuscript\_marked.tex: Line 471}]}\\ A new paragraph has been added to Section~3.4 (highlighted in blue): ``{\itshape \ExecuteMetaData[manuscript_marked.tex]{OverpotentialHER}}'' The references Moraes et al. (IJHE 2026) and Appl. Surf. Sci. 2026, 737, 166886 have been cited.}

\vspace{0.8em}

%--- R1C8 ---
\begin{reviewercomment}[Reviewer 1 --- Comment 8]
  8). Why not give the charge electronic difference for metal-doped $\text{CrCl}_3$ to reveal the electronic interaction between atoms?
\end{reviewercomment}

\response We thank the Reviewer for this question. The charge-density difference and quantitative Bader charge analysis are already presented in the manuscript:
\begin{itemize}[leftmargin=1.5em, itemsep=0.2em]
  \item Figure~4(a--c): Charge-density-difference isosurfaces ($\Delta\rho$) computed as $\Delta\rho = \rho_{\mathrm{TM}/\mathrm{CrCl}_3} - \rho_{\mathrm{CrCl}_3} - \rho_{\mathrm{TM}}$ (Eq.~8), plotted at $\delta=5\times10^{-3}$~$e$/\AA$^{3}$.
  \item Bader charge analysis confirms net electron donation from TM to the substrate: $\text{Co}: +0.77|e|$, $\text{Fe}: +0.94|e|$, $\text{Ni}: +0.59|e|$ (Figure~4(d--f)).
\end{itemize}

\manuscriptchange{\textbf{[Location: \textit{manuscript\_marked.pdf: Page 18, Lines 360--375} \textbar\ \textit{manuscript\_clean.tex: Line 374} \textbar\ \textit{manuscript\_marked.tex: Line 393}]}\\ We have enhanced the discussion surrounding Figure~4 in Section~3.2 (highlighted in blue) to explicitly state the direction of charge transfer and its implications for subsequent H interaction: ``{\itshape \ExecuteMetaData[manuscript_marked.tex]{BaderChargeTransfer}}''}

\vspace{0.8em}

%--- R1C9 ---
\begin{reviewercomment}[Reviewer 1 --- Comment 9]
  9). Furthermore, there are numerous irregularities in the format of the references. It is recommended that the authors strictly adhere to the formatting requirements of this journal, and conduct a systematic review and uniform revision of all the references.
\end{reviewercomment}

\response We thank the Reviewer for identifying these formatting issues. We have conducted a complete, systematic audit of our bibliography (\texttt{references.bib}). All citations have been standardized to strictly comply with the ACS style: complete author listings (eliminating premature \textit{et al.}), full journal title abbreviations, standard volume and page numbers, verified DOIs, correct chemical formula subscripting, and corrected book references (Sholl \& Steckel 2009, Martin 2004). The revised manuscript compiles with zero BibTeX warnings.

\vspace{1.5em}
\noindent\rule{\textwidth}{0.6pt}

\newpage

%========================================================================================
%   RESPONSE TO REVIEWER 2
%========================================================================================
\section{Response to Reviewer \#2}

\noindent\textbf{Recommendation:} Publish after minor revisions noted.

\vspace{0.5em}
\noindent\textbf{General Comments:}\\
\textit{This DFT work systematically investigates Co/Fe/Ni functionalized monolayer $\text{CrCl}_3$ with two coordination modes for HER hydrogen adsorption, combining COHP, AIMD and $d$-band analysis. It clearly reveals that threefold surface Co delivers near-thermoneutral $\Delta G_{\mathrm{ads}}=0.20\text{ eV}$, while sixfold embedding deteriorates activity. However, the paper lacks kinetic barrier calculations for full HER pathway and benchmark comparison with mainstream 2D magnetic catalysts. Spin-orbit coupling is neglected, and only single H coverage is tested, limiting its practical predictive value for real electrolytic environments. I recommend this paper could be accepted if the following questions are addressed properly.}

\vspace{0.5em}
\noindent We thank Reviewer 2 for the positive and constructive evaluation of our study (rating the work in the top 20\% and recommending publication after minor revisions). We address each point below.

\vspace{1em}

%--- R2C1 ---
\begin{reviewercomment}[Reviewer 2 --- Comment 1]
  1. You ignored spin-orbit coupling in all DFT calculations for this magnetic $\text{CrCl}_3$ system; how will including SOC modify TM magnetic moments, orbital hybridization and hydrogen adsorption free energy across adsorbed and embedded configurations?
\end{reviewercomment}

\response We thank the Reviewer for raising this important methodological point.
\begin{enumerate}[leftmargin=1.5em, itemsep=0.2em]
  \item \textbf{Magnitude of SOC in CrCl$_3$:} Unlike $\text{CrI}_3$ where heavy iodine ($Z=53$) induces strong SOC and Ising-like anisotropy, $\text{CrCl}_3$ contains light chlorine ($Z=17$) and is well described by an isotropic Heisenberg model.
  \item \textbf{Effect on TM Moments:} The calculated moments are dominated by intra-atomic exchange; the unquenched orbital moment for $3d$ metals in this coordination is $\sim$0.05--0.12~$\mu_{\mathrm{B}}$, contributing $<$5\% to the total moment.
  \item \textbf{Impact on $\Delta G_{\mathrm{ads}}$:} Covalent TM--Cl and TM--H bond energies are on the scale of 2--7~eV, while SOC splittings are $\sim$10--30~meV. SOC shifts $\Delta G_{\mathrm{ads}}$ by $<$0.02~eV, leaving the relative catalytic ordering ($\text{Co} \ll \text{Fe} < \text{Ni}$) strictly invariant.
\end{enumerate}

\manuscriptchange{\textbf{[Location: \textit{manuscript\_marked.pdf: Page 6, Lines 122--126} \textbar\ \textit{manuscript\_clean.tex: Line 149} \textbar\ \textit{manuscript\_marked.tex: Line 168}]}\\ An explicit paragraph justifying the collinear spin-polarized treatment has been added to Section~2 (Computational Methods, highlighted in blue): ``{\itshape \ExecuteMetaData[manuscript_marked.tex]{SOCStatement}}''}

\vspace{0.8em}

%--- R2C2 ---
\begin{reviewercomment}[Reviewer 2 --- Comment 2]
  2. Only single H coverage was simulated in this study; could you calculate H adsorption at medium/high coverage to explain how adjacent H species compete for active Co sites and shift $\Delta G_{\mathrm{ads}}$ values under real HER conditions?
\end{reviewercomment}

\response We appreciate this insightful comment. Single H coverage (1 H per $2\times 2\times 1$ supercell, corresponding to $1/4\text{ ML}$) was adopted as the standard reference within the CHE framework (N{\o}rskov et al., \textit{J. Phys. Chem. B} 2004, 108, 17886). At higher coverage, adjacent H--H lateral repulsion shifts $\Delta G_{\mathrm{ads}}$ toward more positive values. Since surface Co already has $\Delta G_{\mathrm{ads}} = +0.20\text{ eV}$, increased coverage would maintain the system near the active window. For embedded configurations ($\Delta G_{\mathrm{ads}} > 1.5\text{ eV}$), coverage effects cannot rescue the severely degraded activity. The qualitative ranking is therefore expected to be robust.

\manuscriptchange{\textbf{[Location: \textit{manuscript\_marked.pdf: Pages 24--25, Lines 471--480} \textbar\ \textit{manuscript\_clean.tex: Line 429} \textbar\ \textit{manuscript\_marked.tex: Line 448}]}\\ A dedicated discussion on single-coverage scope and its implications has been added to Section~3.4 (highlighted in blue in the marked manuscript).}

\vspace{0.8em}

%--- R2C3 ---
\begin{reviewercomment}[Reviewer 2 --- Comment 3]
  3. This work only evaluates H adsorption thermodynamics without reaction barriers; can you supply LST/QST transition state energies of Volmer and Tafel steps to fully compare HER kinetics of Co, Fe and Ni catalysts?
\end{reviewercomment}

\response We thank the Reviewer for this comment. The calculation of explicit Volmer/Heyrovsky/Tafel transition-state barriers requires nudged elastic band (NEB) simulations under electrified, solvated interface models, which is computationally prohibitive for the six configurations studied. Within the CHE framework, $\eta = |\Delta G_{\mathrm{ads}}|/e$ serves as a thermodynamic lower bound for the Volmer barrier. Surface-adsorbed Co yields $\eta = 0.20\text{ V}$, confirming kinetic viability.

\manuscriptchange{\textbf{[Location: \textit{manuscript\_marked.pdf: Page 32, Lines 604--606} \textbar\ \textit{manuscript\_clean.tex: Line 475} \textbar\ \textit{manuscript\_marked.tex: Line 494}]}\\ A statement has been added to the Conclusions (Section~4, highlighted in blue): ``{\itshape \ExecuteMetaData[manuscript_marked.tex]{ThermodynamicLimitations}}''}

\vspace{0.8em}

%--- R2C4 ---
\begin{reviewercomment}[Reviewer 2 --- Comment 4]
  4. AIMD simulations only ran for $5\text{ ps}$ at $300\text{ K}$; would longer trajectories or elevated $350/400\text{ K}$ temperatures trigger TM desorption, lattice distortion or irreversible embedding of surface Co and Fe atoms?
\end{reviewercomment}

\response We appreciate this important question. Our 5~ps AIMD trajectories at 300~K establish short-time dynamical stability: Co and Fe remain surface-exposed without thermal desorption or lattice disruption; Ni exhibits spontaneous migration toward the embedded region. At elevated temperatures (350--400~K), thermal activation could facilitate Co/Fe incorporation (thermodynamically favored by $\Delta E_{\mathrm{emb-ads}} = -0.73$ and $-0.86\text{ eV}$). At room temperature, surface-adsorbed Co and Fe remain kinetically trapped, preserving active surface sites.

\manuscriptchange{\textbf{[Location: \textit{manuscript\_marked.pdf: Pages 20--21, Lines 402--412} \textbar\ \textit{manuscript\_clean.tex: Line 393} \textbar\ \textit{manuscript\_marked.tex: Line 412}]}\\ Section~3.2 has been expanded (highlighted in blue): ``{\itshape \ExecuteMetaData[manuscript_marked.tex]{StabilityHierarchy}}''}

\vspace{0.8em}

%--- R2C5 ---
\begin{reviewercomment}[Reviewer 2 --- Comment 5]
  5. You claim coordination engineering tunes activity but never compare $\text{CrCl}_3$ with $\text{CrBr}_3/\text{CrI}_3$; how do heavier halide ligands alter $\text{TM--Cl}$ bonding strength and optimize H adsorption for the same threefold Co active sites?
\end{reviewercomment}

\response We thank the Reviewer for this comparative suggestion. $\text{CrCl}_3$ was chosen because its higher chemical stability and harder Cl coordination maximize active-site localization. Extension to heavier CrX$_3$ homologs would require recalibrated pseudopotentials, U parameters, and COHP projectors. The expected qualitative trend (larger lattice, more diffuse p-states, modified crystal-field splitting) may reduce ICOHP magnitudes and increase site exposure.

\manuscriptchange{\textbf{[Location: \textit{manuscript\_marked.pdf: Page 32, Lines 606--610} \textbar\ \textit{manuscript\_clean.tex: Line 475} \textbar\ \textit{manuscript\_marked.tex: Line 494}]}\\ A comparative statement has been added to the Conclusions (Section~4, highlighted in blue): ``{\itshape \ExecuteMetaData[manuscript_marked.tex]{BroaderTransferability}}''}

\vspace{0.8em}

%--- R2C6 ---
\begin{reviewercomment}[Reviewer 2 --- Comment 6]
  6. The $d$-band center shows weak linear correlation with $\Delta G_{\mathrm{ads}}$; what combined electronic descriptor integrating spin splitting and ICOHP values can better predict hydrogen binding strength on magnetic TM centers?
\end{reviewercomment}

\response We thank the Reviewer for this perceptive question. As discussed in Section~3.5 and Figure~8, the occupied $d$-band center ($\varepsilon_d^{\mathrm{occ}}$) shows weak $R^2$ ($0.362$ for surface-adsorbed, $0.770$ for embedded). This failure arises because a spin-summed first moment cannot capture spin-channel polarization, crystal-field splitting, or structural deformation penalties. A combined descriptor integrating spin-resolved $d$-band centers and cumulative TM--Cl ICOHP would be more predictive, but with only three metals per coordination family, fitting such a multi-parameter descriptor is not statistically meaningful.

\manuscriptchange{\textbf{[Location: \textit{manuscript\_marked.pdf: Pages 29--30, Lines 542--569} \textbar\ \textit{manuscript\_clean.tex: Line 456} \textbar\ \textit{manuscript\_marked.tex: Line 475}]}\\ An expanded paragraph discussing the combined descriptor and its limitations has been added to Section~3.5 (highlighted in blue in the marked manuscript).}

\vspace{0.8em}

%--- R2C7 ---
\begin{reviewercomment}[Reviewer 2 --- Comment 7]
  7. No direct benchmark against classic HER catalysts like Pt or single-atom M-N-C was provided under identical DFT parameters; can you compute their $\Delta G_{\mathrm{ads}}$ to quantify $\text{Co-CrCl}_3$'s competitive advantages?
\end{reviewercomment}

\response We appreciate this suggestion. The well-established literature benchmarks are:
\begin{itemize}[leftmargin=1.5em, itemsep=0.2em]
  \item Pt(111): $\Delta G_{\mathrm{ads}} \approx -0.09\text{ eV}$ (near-thermoneutral, slightly overbinding);
  \item Co-N$_4$-C (M-N-C): $\Delta G_{\mathrm{ads}} \approx +0.10$ to $+0.15\text{ eV}$;
  \item Surface-adsorbed Co/CrCl$_3$ (this work): $\Delta G_{\mathrm{ads}} = 0.20\text{ eV}$ ($\eta = 0.20\text{ V}$).
\end{itemize}
Recalculating Pt and M-N-C under identical DFT parameters would constitute a separate benchmarking study; however, the literature values provide the relevant context.

\manuscriptchange{\textbf{[Location: \textit{manuscript\_marked.pdf: Page 28, Lines 527--536} \textbar\ \textit{manuscript\_clean.tex: Line 452} \textbar\ \textit{manuscript\_marked.tex: Line 471}]}\\ A benchmark comparison paragraph has been added to Section~3.4 (highlighted in blue in the marked manuscript), citing N{\o}rskov 2005 and Feng 2021.}

\vspace{0.8em}

%--- R2C8 ---
\begin{reviewercomment}[Reviewer 2 --- Comment 8]
  8. Monolayer $\text{CrCl}_3$ suffers from high raw material and manufacturing costs, which restricts its large-scale commercial hydrogen production, even with decent theoretical HER activity. Low-cost carbon or metal oxide catalysts are more suitable for industrial mass hydrogen generation.
\end{reviewercomment}

\response We thank the Reviewer for this practical observation. We agree regarding industrial scaling. The primary aim of this study is to elucidate the fundamental physics of spin--coordination coupling in magnetic 2D van der Waals systems, not to propose a mass-market replacement for carbon or oxide electrocatalysts. The coordination engineering principles identified (threefold vs. sixfold coordination; active-site exposure as a design rule) are transferable to more cost-effective systems, including TM-doped carbon materials, MXenes, and 2D chalcogenides.

\manuscriptchange{\textbf{[Location: \textit{manuscript\_marked.pdf: Page 32, Lines 610--612} \textbar\ \textit{manuscript\_clean.tex: Line 475} \textbar\ \textit{manuscript\_marked.tex: Line 494}]}\\ A statement about broader transferability has been added to the Conclusions (highlighted in blue).}

\vspace{0.8em}

%--- R2C9 ---
\begin{reviewercomment}[Reviewer 2 --- Comment 9]
  9. It is suggested that you cite the provided references [DOI: 10.1016/j.ijhydene.2020.11.102; https://doi.org/10.1016/j.ijhydene.2022.08.200; DOI: 10.1016/j.jallcom.2026.187619] to supplement and deepen the background description in the introduction.
\end{reviewercomment}

\response We thank the Reviewer for these recommendations. We have reviewed all three papers:
\begin{itemize}[leftmargin=1.5em, itemsep=0.2em]
  \item Feng et al., \textit{IJHE} 2021, 46, 5378: Cu embedded in graphdiyne for HER/CO$_2$RR using the CHE framework -- directly relevant to our finding that coordination controls $\Delta G_{\mathrm{ads}}$.
  \item Feng et al., \textit{IJHE} 2022, 47, 36294: Lattice strain in biphenylene for ORR -- broadly consistent with our coordination engineering paradigm.
  \item Feng et al., \textit{J. Alloys Compd.} 2026, 1062, 187619: Ni$_3$S$_2$ for HMF oxidation -- cited in the context of TM electrocatalysts for clean energy.
\end{itemize}

\manuscriptchange{\textbf{[Location: \textit{manuscript\_marked.pdf: Page 3, Lines 39--40} \textbar\ \textit{manuscript\_clean.tex: Line 130} \textbar\ \textit{manuscript\_marked.tex: Line 149}]}\\ All three references have been cited in the Introduction (highlighted in blue): \texttt{\textbackslash cite\{Feng2021IJHE\}}, \texttt{\textbackslash cite\{Feng2022IJHE\}}, \texttt{\textbackslash cite\{Feng2026JAC\}}.}

\vspace{1.5em}
\noindent\rule{\textwidth}{0.6pt}

\newpage

%========================================================================================
%   RESPONSE TO REVIEWER 3
%========================================================================================
\section{Response to Reviewer \#3}

\noindent\textbf{Recommendation:} Not appropriate for ACS Applied Energy Materials.

\vspace{0.5em}
\noindent\textbf{General Comments:}\\
\textit{In this manuscript, spin-polarized DFT calculation and ab initio MD simulations are adopted to investigate the transition metal functionalization and the effects of site coordination on hydrogen adsorption on monolayer $\text{CrCl}_3$. It was found that surface-adsorbed metals introduce localized, spin-polarized $3d$ states that drastically enhance HER, and the strong substrate binding does not guarantee optimal catalytic activity. Transitioning from threefold surface sites to sixfold embedded configurations systematically suppresses H binding, and surface-exposed sites optimize the balance between catalyst stability and active-site reactivity. The topic and results are interesting, however, the it seems too simple to evaluate the HER properties only according to the hydrogen adsorption free energy. So I cannot recommend it to be considered for publication in its present version. The main comments are summarized as:}

\vspace{0.5em}
\noindent We thank Reviewer 3 for the constructive and technically detailed comments. Regarding the general concern that ``it seems too simple to evaluate HER properties only according to $\Delta G_{\mathrm{ads}}$'': we respectfully note that the computational hydrogen electrode (CHE) framework using $\Delta G_{\mathrm{ads}}$ as the primary HER descriptor is the established and widely accepted methodology in the field (N{\o}rskov et al., \textit{PCCP} 2004; \textit{JPCB} 2004; 2005; Moraes et al. 2026; Feng et al. 2021), used by hundreds of DFT-based catalyst screening studies.

\vspace{1em}

%--- R3C1 ---
\begin{reviewercomment}[Reviewer 3 --- Comment 1]
  (1) The hydrogen adsorption free energy calculated in this work is $1.82\text{ eV}$, which is much larger than that ($0.13\text{ eV}$) reported in reference. Why? The authors should check and modify the data accurately. Otherwise, the following results are questionable.
\end{reviewercomment}

\response We thank the Reviewer for highlighting this important point. We emphasize that both values -- our $1.82\text{ eV}$ and the literature $0.13\text{ eV}$ (Zala et al., 2024) -- are physically correct, but correspond to \textbf{fundamentally different adsorption geometries}:
\begin{itemize}[leftmargin=1.5em, itemsep=0.2em]
  \item \textbf{This Work (Site S$_3$ --- Direct Top-Cr Chemisorption):} H penetrates the Cl surface layer and forms a direct H--Cr bond ($d_{\mathrm{H-Cr}} = 1.55\text{ \AA}$). This incurs a substantial structural/electrostatic penalty, yielding $\Delta E_{\mathrm{ads}} = +1.58\text{ eV}$ and $\Delta G_{\mathrm{ads}} = 1.82\text{ eV}$ after the standard $+0.24\text{ eV}$ ZPE-entropy correction.
  \item \textbf{Literature (Zala et al., 2024 --- On-Cl Configuration):} H interacts weakly with the outermost Cl atoms without reaching Cr.
  \item \textbf{Sites S$_1$ and S$_2$:} H placed at top-Cl (S$_1$) or hollow (S$_2$) desorbs ($d > 3.42\text{ \AA}$), confirming no stable chemisorbed state at these sites under our protocol.
\end{itemize}
Both values demonstrate that pristine CrCl$_3$ has poor intrinsic HER activity, motivating TM functionalization.

\manuscriptchange{\textbf{[Location: \textit{manuscript\_marked.pdf: Pages 15--16, Lines 315--324} \textbar\ \textit{manuscript\_clean.tex: Line 355} \textbar\ \textit{manuscript\_marked.tex: Line 374}]}\\ An explicit comparison has been expanded in Section~3.1 (highlighted in blue): ``{\itshape \ExecuteMetaData[manuscript_marked.tex]{PristineGadsComparison}}''}

\vspace{0.8em}

%--- R3C2 ---
\begin{reviewercomment}[Reviewer 3 --- Comment 2]
  (2) The atom-projected moments of surface-adsorbed Co, Fe, and Ni are $1.99$, $3.34$, and $-0.76\,\mu_{\mathrm{B}}$, respectively, Co and Fe are polarized parallel to the ferromagnetic host, whereas Ni is antiparallel. Why? How can we understand it from the intrinsic physics?
\end{reviewercomment}

\response We thank the Reviewer for this physically deep question. The magnetic behavior of each TM dopant can be understood from three competing factors: (i)~$d$-electron count and Hund's first rule, (ii)~crystal-field splitting at the threefold Cl coordination site, and (iii)~exchange coupling to the ferromagnetic Cr sublattice:
\begin{enumerate}[leftmargin=1.5em, itemsep=0.2em]
  \item \textbf{Iron ($d^6$, $+3.34\,\mu_{\mathrm{B}}$, parallel):} High-spin configuration ($S \approx 2$); Cr--Cl--Fe super-exchange following Goodenough--Kanamori rules stabilizes ferromagnetic alignment.
  \item \textbf{Cobalt ($d^7$, $+1.99\,\mu_{\mathrm{B}}$, parallel):} Intermediate-spin state; parallel coupling via Cr--Cl--Co super-exchange (half-filled $\to$ partially-filled overlap).
  \item \textbf{Nickel ($d^8$, $-0.76\,\mu_{\mathrm{B}}$, antiparallel):} Nearly-full majority-spin channel blocks majority-spin hopping (Pauli exclusion). The dominant kinetic exchange is via minority-spin Ni states coupling to Cr, yielding antiferromagnetic super-exchange.
\end{enumerate}
This is consistent with the TM-doping literature for CrX$_3$ systems, where $d^8$/$d^9$ metals tend toward antiparallel polarization.

\manuscriptchange{\textbf{[Location: \textit{manuscript\_marked.pdf: Pages 22--24, Lines 430--470} \textbar\ \textit{manuscript\_clean.tex: Lines 412--428} \textbar\ \textit{manuscript\_marked.tex: Lines 431--447}]}\\ A comprehensive physical rationalization has been added to Section~3.2 (highlighted in blue in the marked manuscript):}
\begin{addedtext}
  Co, Fe, and Ni carry $d^7$, $d^6$, and $d^8$ electron configurations, respectively, placing them near the thermoneutral region of the HER volcano plot where $\Delta G_{\mathrm{ads}}$ can be tuned toward zero by controlling the local coordination environment. Earlier $3d$ metals (e.g., Ti, V) tend to overbind H, while late metals with filled $d$-shells (e.g., Cu, Zn) offer limited hybridization. Co, Fe, and Ni also carry significant magnetic moments ($1$--$4~\mu_{\mathrm{B}}$) compatible with the Cr$^{3+}$ ($d^3$) Heisenberg ferromagnet ground state of CrCl$_3$, enabling spin--coordination--catalysis coupling. They are furthermore among the most commonly reported TM functionalization agents for CrCl$_3$ and related chromium trihalides in both theoretical and experimental studies, establishing their experimental realizability.
\end{addedtext}

\vspace{0.8em}

%--- R3C3 ---
\begin{reviewercomment}[Reviewer 3 --- Comment 3]
  (3) It is stated that embedding reorganizes the TM-$3d$ manifold through the change from threefold to distorted sixfold coordination. How did the authors obtain such a conclusion from Figure 5(g--l)?
\end{reviewercomment}

\response We thank the Reviewer for requesting greater clarity. The conclusion is based on three observable spectral changes visible in Figure~5(g--l) compared to Figure~5(a--f):
\begin{itemize}[leftmargin=1.5em, itemsep=0.2em]
  \item \textbf{Peak Redistribution:} Surface-adsorbed Co shows a single broad TM-$3d$ feature near $E_{\mathrm{F}}$; embedded Co shows multiple narrow peaks spread over a wider energy window ($\sim$$-5$ to $-1\text{ eV}$), indicating enhanced crystal-field splitting.
  \item \textbf{$E_{\mathrm{F}}$ Density Reduction:} Embedded Co and Ni show fewer TM-$3d$ states at the Fermi level, explaining reduced orbital availability for H-$1s$ hybridization (consistent with weaker $E_{\mathrm{int}}^{\mathrm{H}}$: Co $-0.94$ vs. $-2.46\text{ eV}$).
  \item \textbf{Enhanced TM-$3d$/Cl-$3p$ Mixing:} In panels (g--l), TM-$3d$ states show stronger spectral overlap with the Cl-$3p$ band ($\sim$$-6$ to $-2\text{ eV}$), confirmed by the cumulative ICOHP values jumping from $\sim$$-6.7\text{ eV}$ (threefold) to $\sim$$-11.2\text{ eV}$ (sixfold) for Co.
\end{itemize}

\manuscriptchange{\textbf{[Location: \textit{manuscript\_marked.pdf: Page 21, Lines 401--410} \textbar\ \textit{manuscript\_clean.tex: Line 401} \textbar\ \textit{manuscript\_marked.tex: Line 420}]}\\ The paragraph discussing Figure~5(g--l) has been substantially expanded in Section~3.3 (highlighted in blue) to provide this explicit reading guide.}

\vspace{1.5em}
\noindent\rule{\textwidth}{0.6pt}

\newpage

%========================================================================================
%   RESPONSE TO REVIEWER 4 & 6
%========================================================================================
\section{Response to Reviewer \#4 (and Reviewer \#6)}

\noindent\textbf{Recommendation:} Major revisions needed as noted.

\vspace{0.5em}
\noindent\textbf{General Comments:}\\
\textit{This manuscript presents a spin-polarized DFT and ab initio molecular dynamics (AIMD) investigation of Co-, Fe-, and Ni-functionalized monolayer $\text{CrCl}_3$\ldots Nevertheless, it seems that the manuscript requires major revision before it can be considered for publication.}

\vspace{0.5em}
\noindent We sincerely thank Reviewer 4 for the exceptionally thorough and rigorous review. The extensive comments identify genuine opportunities to strengthen the work. Below, we address all 20 points individually.

\vspace{1em}

%--- R4C1 ---
\begin{reviewercomment}[Reviewer 4 --- Comment 1]
  1. Plain PBE has been used throughout for a system in which correlation matters. $\text{CrCl}_3$ is a correlated $3d$ magnetic insulator, and the reported PBE gap of $1.50\text{ eV}$ is well below the experimentally reported optical gap of roughly $3\text{ eV}$. Self-interaction error in semi-local GGA is known to misplace $3d$ levels\ldots Why was DFT+U (or a hybrid functional) not employed, at least as a benchmark?
\end{reviewercomment}

\response We thank the Reviewer for this important methodological remark. While PBE underestimates the optical gap ($1.50$ vs. $\sim$3~eV), this limitation does not affect the qualitative relative trends, which are governed by TM-$3d$ valence-electron count and TM--Cl/TM--H orbital hybridization, not host excitation energies. Standard benchmark studies on CrCl$_3$ (Zala2024, Ipaves2025, Liao2022) routinely employ the same approach. The $\Delta G_{\mathrm{ads}}$ ordering (Co $<$ Fe $<$ Ni) is robustly supported by $d$-electron count and ICOHP arguments independent of $U$.

\manuscriptchange{\textbf{[Location: \textit{manuscript\_marked.pdf: Page 6, Lines 127--135} \textbar\ \textit{manuscript\_clean.tex: Line 152} \textbar\ \textit{manuscript\_marked.tex: Line 171}]}\\ A new paragraph has been added to Section~2 (highlighted in blue): ``{\itshape \ExecuteMetaData[manuscript_marked.tex]{PBEHubbardU}}''}

\vspace{0.8em}

%--- R4C2 ---
\begin{reviewercomment}[Reviewer 4 --- Comment 2]
  2. No van der Waals correction is mentioned anywhere in Computational Methods. Dispersion is not negligible for a metal adatom on a chloride surface, and it is decisive for the pristine reference case in Figure 2(b), where H relaxes to $3.42\text{ \AA}$\ldots Please state explicitly that no dispersion scheme was applied and justify that choice or add DFT-D3 (or equivalent) checks.
\end{reviewercomment}

\response We thank the Reviewer for identifying this important methodological point. In the main manuscript, plain PBE was employed to maintain direct consistency with established benchmarks on 2D chromium halides (Zala2024, Luo2020, Mastrippolito2022). For the functionalized systems, the dominant TM--CrCl$_3$ and TM--H interactions are strongly covalent, as confirmed by cumulative TM--Cl ICOHP values ranging from $-6.2$ to $-11.4\text{ eV}$ and short TM--H bond lengths (1.43--1.57~\AA).

To definitively address the Reviewer's request, we have performed explicit \textbf{DFT-D3 (Grimme zero-damping, IVDW=11)} benchmark calculations across the full multi-scale matrix of supercells ($1\times 1$, $2\times 2$, and $3\times 3$) and all three adsorption sites ($S_1$, $S_2$, and $S_3$) for pristine monolayer $\text{CrCl}_3$. The calculated results are consolidated in the table below:

\begin{center}
  \small
  \resizebox{\textwidth}{!}{%
  \begin{tabular}{lcccccc}
    \toprule
    Supercell & Coverage $\theta$ & Adsorption Site & Pure PBE $\Delta G_{\mathrm{ads}}$ (eV) & PBE+D3 $\Delta G_{\mathrm{ads}}$ (eV) & $\Delta(\text{D3}-\text{PBE})$ (eV) & Relative Stability \\
    \midrule
    $1\times 1$ & $1.000$ & $S_1$ (Top-Cl) & 1.918 & 1.781 & $-0.137$ & Thermodynamically favored \\
                &         & $S_2$ (Hollow) & 2.727 & 2.617 & $-0.110$ & Strongly disfavored \\
                &         & $S_3$ (Top-Cr) & 1.876 & 1.895 & $+0.019$ & Apical chemisorbed \\
    \midrule
    $2\times 2$ (paper) & $0.250$ & $S_1$ (Top-Cl) & 1.819 & 1.753 & $-0.066$ & Thermodynamically favored \\
                &         & $S_2$ (Hollow) & 2.711 & 2.599 & $-0.112$ & Strongly disfavored \\
                &         & $S_3$ (Top-Cr) & 1.864 & 1.864 & $-0.0005$ & Apical chemisorbed \\
    \midrule
    $3\times 3$ (dilute)& $0.111$ & $S_1$ (Top-Cl) & 1.506 & 1.537 & $+0.031$ & Thermodynamically favored \\
                &         & $S_2$ (Hollow) & 2.710 & 2.563 & $-0.147$ & Strongly disfavored \\
                &         & $S_3$ (Top-Cr) & 1.840 & 1.869 & $+0.029$ & Apical chemisorbed \\
    \bottomrule
  \end{tabular}%
  }
\end{center}

These benchmarks demonstrate:
\begin{enumerate}[leftmargin=1.5em, itemsep=0.2em]
  \item \textbf{Chemisorption Invariance at Metal Center:} At the chemisorbed $S_3$ (top-Cr) site, which directly parallels the apical TM--H chemisorption bond, the dispersion shift is negligible: $\Delta(\text{D3}-\text{PBE}) \le 0.03\text{ eV}$ across all supercells, and exactly $0.00\text{ eV}$ for the $2\times 2$ paper benchmark.
  \item \textbf{Preservation of Site Stability Hierarchy:} The stability sequence ($S_1 < S_3 \ll S_2$) is rigorously preserved with and without dispersion across all three supercell dimensions.
  \item \textbf{Robustness of Inactive Pristine State:} For all configurations and cell sizes, $\Delta G_{\mathrm{ads}} \ge 1.51\text{ eV}$, confirming that pristine $\text{CrCl}_3$ remains strongly endergonic and catalytically inactive regardless of dispersion inclusion.
\end{enumerate}

\manuscriptchange{\textbf{[Location: \textit{manuscript\_marked.pdf: Pages 6--7, Lines 136--144} \textbar\ \textit{manuscript\_clean.tex: Line 154} \textbar\ \textit{manuscript\_marked.tex: Line 173}]}\\ An explicit statement incorporating these DFT-D3 benchmarks has been added to Section~2 (highlighted in blue): ``{\itshape \ExecuteMetaData[manuscript_marked.tex]{DispersionCorrection}}''}

\vspace{0.8em}

%--- R4C3 ---
\begin{reviewercomment}[Reviewer 4 --- Comment 3]
  3. The binding energies are referenced to isolated spin-polarized transition-metal atoms\ldots The PBE cohesive energies of bulk Co, Fe, and Ni ($\sim$4.9--5.5~eV/atom) exceed the reported $|E_{\mathrm{bind}}|$ for every adsorbed configuration\ldots Please quantify this by reporting $E_{\mathrm{bind}} - E_{\mathrm{coh}}$.
\end{reviewercomment}

\response We thank the Reviewer for this valid and important point. The comparison is:

\begin{center}
  \small
  \begin{tabular}{lcccl}
    \toprule
    System      & $|E_{\mathrm{bind}}|$ (eV) & $E_{\mathrm{coh,bulk}}$ (eV) & $\Delta$ (eV) & Stability                \\
    \midrule
    Co-adsorbed & 2.82                       & 4.39                         & +1.57         & clustering preferred     \\
    Fe-adsorbed & 2.60                       & 4.28                         & +1.68         & clustering preferred     \\
    Ni-adsorbed & 3.38                       & 4.44                         & +1.06         & clustering preferred     \\
    Co-embedded & 3.55                       & 4.39                         & +0.84         & clustering preferred     \\
    Fe-embedded & 3.36                       & 4.28                         & +0.92         & clustering preferred     \\
    Ni-embedded & 4.95                       & 4.44                         & $-0.51$       & \textbf{stable vs. bulk} \\
    \bottomrule
  \end{tabular}
\end{center}
Only embedded Ni is thermodynamically stable against bulk formation. Adsorbed configurations are kinetically stabilized by the energy barrier for diffusion through the constrained CrCl$_3$ hollow site.

\manuscriptchange{\textbf{[Location: \textit{manuscript\_marked.pdf: Page 17, Lines 348--357} \textbar\ \textit{manuscript\_clean.tex: Line 372} \textbar\ \textit{manuscript\_marked.tex: Line 391}]}\\ A new paragraph has been added to Section~3.2 (highlighted in blue): ``{\itshape \ExecuteMetaData[manuscript_marked.tex]{BulkCohesiveComparison}}''}

\vspace{0.8em}

%--- R4C4 ---
\begin{reviewercomment}[Reviewer 4 --- Comment 4]
  4. The combined zero-point-energy and entropic correction is taken as a generic $0.24\text{ eV}$ for all systems. Given that the $\text{TM--H}$ distances span $1.43\text{--}1.57\text{ \AA}$\ldots Please compute the vibrational frequencies of the adsorbed H for each configuration and use system-specific $\Delta E_{\mathrm{ZPE}} - T\Delta S$ values.
\end{reviewercomment}

\response We thank the Reviewer for this legitimate methodological point. The $+0.24\text{ eV}$ correction is the standard CHE benchmark (N{\o}rskov et al., \textit{JPCB} 2004). For TM--H distances spanning 1.43--1.57~\AA, the H stretching mode varies by $\sim$100--200~cm$^{-1}$, contributing at most $\sim$10--20~meV to the correction -- well below the $\Delta G_{\mathrm{ads}}$ differences between metals (0.2--0.7~eV). System-specific DFPT vibrational frequencies are identified as a future improvement.

\manuscriptchange{\textbf{[Location: \textit{manuscript\_marked.pdf: Page 8, Lines 174--177} \textbar\ \textit{manuscript\_clean.tex: Line 188} \textbar\ \textit{manuscript\_marked.tex: Line 207}]}\\ A statement acknowledging the generic correction and its known uncertainty has been added to Section~2 (Computational Methods, highlighted in blue).}

\vspace{0.8em}

%--- R4C5 ---
\begin{reviewercomment}[Reviewer 4 --- Comment 5]
  5. The electrical conductivity of the support is not addressed\ldots Do the TM-derived in-gap states shown in Figure 5 form a continuous conduction channel at realistic dopant concentrations, or are they localized?
\end{reviewercomment}

\response We thank the Reviewer for raising this important applied question. The TM-derived in-gap states (Figure~5) are weakly dispersive and localized around the functionalizing center; they do not form a percolating conduction channel at the dilute dopant concentration modeled. Practical implementation requires a conducting support electrode (carbon cloth, graphene, or glassy carbon), consistent with how 2D-semiconductor-based HER catalysts are conventionally deployed.

\manuscriptchange{\textbf{[Location: \textit{manuscript\_marked.pdf: Pages 24--25, Lines 471--480} \textbar\ \textit{manuscript\_clean.tex: Line 429} \textbar\ \textit{manuscript\_marked.tex: Line 448}]}\\ A new paragraph has been added to Section~3.3 (highlighted in blue): ``{\itshape \ExecuteMetaData[manuscript_marked.tex]{ConductivityBandGap}}''}

\vspace{0.8em}

%--- R4C6 ---
\begin{reviewercomment}[Reviewer 4 --- Comment 6]
  6. It appears that only one collinear magnetic solution was converged per system. The antiparallel Ni moment ($-0.76\,\mu_{\mathrm{B}}$) is unusual\ldots Were parallel and antiparallel alignments both converged and compared energetically?
\end{reviewercomment}

\response We appreciate this important point. In our calculations, all initial TM moments were set \textbf{parallel} to the Cr sublattice (MAGMOM = $8\times3.0$ (Cr) $24\times0.0$ (Cl) $1\times3.0$ (TM)). The Ni antiparallel state converged \textbf{spontaneously} from this parallel initialization, which is a strong indicator that the antiparallel state is the self-consistent ground state. This is physically expected based on Goodenough--Kanamori--Anderson super-exchange rules for $d^8$ Ni coupled to $d^3$ Cr through Cl ligands (see Response to Reviewer 3, Comment 2).

\manuscriptchange{\textbf{[Location: \textit{manuscript\_marked.pdf: Page 18, Lines 383--387} \textbar\ \textit{manuscript\_clean.tex: Line 385} \textbar\ \textit{manuscript\_marked.tex: Line 404}]}\\ A statement specifying the spontaneous convergence to the antiparallel Ni state has been added to Section~2 (Computational Methods, highlighted in blue).}

\vspace{0.8em}

%--- R4C7 ---
\begin{reviewercomment}[Reviewer 4 --- Comment 7]
  7. The AIMD computational details is not properly specified. The time step, the k-point sampling used during the dynamics, the Nosé--Hoover coupling parameter, and the equilibration procedure are not given anywhere\ldots
\end{reviewercomment}

\response We thank the Reviewer for identifying this gap. The complete AIMD parameters from our INCAR files are:

\begin{center}
  \small
  \begin{tabular}{ll}
    \toprule
    Parameter                      & Value                                     \\
    \midrule
    Time step (POTIM)              & 2 fs                                      \\
    Number of steps (NSW)          & 2500 (= 5 ps total)                       \\
    Temperature (TEBEG/TEEND)      & 300 K / 300 K                             \\
    Thermostat                     & Nosé--Hoover (MDALGO = 2)                 \\
    Coupling (SMASS)               & 0 (default)                               \\
    $k$-point sampling             & $\Gamma$-point only (1$\times$1$\times$1) \\
    Electronic convergence (EDIFF) & $10^{-5}$ eV                              \\
    Equilibration                  & First 0.5 ps (250 steps)                  \\
    \bottomrule
  \end{tabular}
\end{center}

\manuscriptchange{\textbf{[Location: \textit{manuscript\_marked.pdf: Page 12, Lines 255--263} \textbar\ \textit{manuscript\_clean.tex: Line 325} \textbar\ \textit{manuscript\_marked.tex: Line 344}]}\\ All AIMD computational details have been added to Section~2 (highlighted in blue): ``{\itshape \ExecuteMetaData[manuscript_marked.tex]{AIMDParameters}}''}

\vspace{0.8em}

%--- R4C8 ---
\begin{reviewercomment}[Reviewer 4 --- Comment 8]
  8. In Figure S3 the total energy varies by roughly $4\text{ eV}$ over the course of the trajectories and shows numerous sharp downward spikes. Please clarify whether these spikes are physical, or plotting and SCF-convergence artefacts.
\end{reviewercomment}

\response We thank the Reviewer for this observation. The sharp downward spikes are SCF convergence artifacts, not physical events. During NVT-AIMD, the OSZICAR records the total energy after each ionic step. Normal thermal fluctuations of $\sim$1--3~eV over 5~ps at 300~K are expected for a 33-atom supercell. The isolated downward spikes occur when the electronic SCF converges to a lower-energy metastable state at specific ionic configurations. Energy conservation is confirmed by the temperature profiles (Figure S3 lower panels), which fluctuate appropriately around 300~K.

\manuscriptchange{\textbf{[Location: \textit{manuscript\_marked.pdf: Page 20, Lines 391--392 \& SI Fig. S3} \textbar\ \textit{manuscript\_clean.tex: Line 391} \textbar\ \textit{manuscript\_marked.tex: Line 410}]}\\ A clarifying note has been added to the Supporting Information caption for Figure S3 (highlighted in blue).}

\vspace{0.8em}

%--- R4C9 ---
\begin{reviewercomment}[Reviewer 4 --- Comment 9]
  9. Details and sign conventions of the Bader analysis need clarification. Define the direction of the reported charge transfer\ldots the color-bar label in Figure 4(d--f) reads $\Delta|Q_e|$ but is used with negative values, which is contradictory notation; use $\Delta Q\text{ }(e)$.
\end{reviewercomment}

\response We thank the Reviewer for identifying this notation inconsistency. The TM adatoms \textbf{donate} electrons to the CrCl$_3$ substrate. The corrected values are: Co donates $0.77|e|$, Fe donates $0.94|e|$, Ni donates $0.59|e|$.

\manuscriptchange{\textbf{[Location: \textit{manuscript\_marked.pdf: Page 18, Lines 361--363} \textbar\ \textit{manuscript\_clean.tex: Line 374} \textbar\ \textit{manuscript\_marked.tex: Line 393}]}\\ The text in Section~3.2 has been corrected (highlighted in blue): ``{\itshape \ExecuteMetaData[manuscript_marked.tex]{BaderChargeTransfer}}'' The Figure~4(d--f) color-bar label has been corrected from ``$\Delta|Q_e|$'' to ``$\Delta Q\,(e)$'' with a signed convention.}

\vspace{0.8em}

%--- R4C10 ---
\begin{reviewercomment}[Reviewer 4 --- Comment 10]
  10. The electronic interpretation of the Ni case deserves more than the observation that the trends do not correlate. Ni shows the smallest charge transfer, the strongest binding, and an antiparallel local moment. Please provide clearer rationalization\ldots
\end{reviewercomment}

\response We thank the Reviewer for pressing for greater physical depth. The Ni anomaly is rationalized by its $d^8$ configuration:
\begin{itemize}[leftmargin=1.5em, itemsep=0.2em]
  \item The nearly-full majority-spin channel leaves little capacity for charge donation, explaining the \textbf{smallest charge transfer} (0.59$|e|$ vs. 0.77/0.94 for Co/Fe).
  \item The \textbf{strongest binding} despite smallest charge transfer arises from shorter Ni--Cl bonds (2.16~\AA) providing strong Pauli-repulsion-free orbital overlap and additional covalent bonding from filled minority-spin states.
  \item The \textbf{antiparallel moment} ($-0.76\,\mu_{\mathrm{B}}$) reflects antiferromagnetic super-exchange of the nearly-full $d^8$ to the Cr host through Cl.
\end{itemize}

\manuscriptchange{\textbf{[Location: \textit{manuscript\_marked.pdf: Page 24, Lines 425--428} \textbar\ \textit{manuscript\_clean.tex: Line 427} \textbar\ \textit{manuscript\_marked.tex: Line 446}]}\\ This analysis has been added to Section~3.2 (highlighted in blue in the marked manuscript):}
\begin{addedtext}
  The local magnetic moment leads to the same conclusion (Figure~S5). Co provides the most favorable surface adsorption despite having neither the largest TM moment nor antiparallel polarization, and embedding changes the Ni moment only moderately while increasing its $\Delta G_{\mathrm{ads}}$ from 0.88 to 3.62~eV. Neither the magnitude nor the sign of the H-free local moment alone explains the adsorption trend.
\end{addedtext}

\vspace{0.8em}

%--- R4C11 ---
\begin{reviewercomment}[Reviewer 4 --- Comment 11]
  11. The manuscript contains no tables at all\ldots Please add one or two consolidated tables summarizing all quantities for the seven systems.
\end{reviewercomment}

\response We thank the Reviewer for this practical suggestion. We fully agree that consolidated tables significantly improve readability. We have added summary tables to the Results and Discussion section, consolidating all key quantities ($E_{\mathrm{bind}}$, $\Delta E_{\mathrm{emb-ads}}$, $\mu_{\mathrm{TM}}$, $Q_{\mathrm{Bader}}$, ICOHP$_{\mathrm{TM-Cl}}$, $\varepsilon_d^{\mathrm{occ}}$, $\Delta G_{\mathrm{ads}}$) for the 7 systems (pristine + 6 functionalized).

\manuscriptchange{\textbf{[Location: \textit{manuscript\_marked.pdf: Page 13, Table 1 \& Page 27, Table 2} \textbar\ \textit{manuscript\_clean.tex: Lines 333, 441} \textbar\ \textit{manuscript\_marked.tex: Lines 352, 460}]}\\ New consolidated tables have been added to the Results section (highlighted in blue in the marked manuscript).}

\vspace{0.8em}

%--- R4C12 ---
\begin{reviewercomment}[Reviewer 4 --- Comment 12]
  12. Notation is not uniform. The manuscript uses $\Delta G_{\mathrm{ads}}$ in the text and in Figures 7 and S5, $\Delta G_{\mathrm{H}}$ in Figure 8, and $G_{\mathrm{ads}}$ in the bar-chart legends of Figure 7\ldots Please unify the notation throughout.
\end{reviewercomment}

\response We thank the Reviewer for identifying these inconsistencies. All notation has been rigorously unified to $\Delta G_{\mathrm{ads}}$ throughout the manuscript, figures, and Supporting Information. Specifically: Figure~8 axis label changed from $\Delta G_{\mathrm{H}}$ to $\Delta G_{\mathrm{ads}}$; Figure~7 legend changed from $G_{\mathrm{ads}}$ to $\Delta G_{\mathrm{ads}}$; H$^*$ notation is used consistently and defined on first use.

\manuscriptchange{\textbf{[Location: \textit{manuscript\_marked.pdf: Throughout manuscript} \textbar\ \textit{manuscript\_clean.tex: Multiple sections} \textbar\ \textit{manuscript\_marked.tex: Multiple sections}]}\\ Notation corrections are highlighted in blue where substantive; noted as editorial corrections elsewhere.}

\vspace{0.8em}

%--- R4C13 ---
\begin{reviewercomment}[Reviewer 4 --- Comment 13]
  13. The reference list needs substantial correction; it does not follow ACS Applied Energy Materials format\ldots
\end{reviewercomment}

\response We thank the Reviewer for the detailed identification of every bibliographic issue. Every reference has been corrected in \texttt{references.bib}: Ref.~46 (Sholl \& Steckel textbook) restored; Ref.~47 (Martin) corrected; Ref.~31 duplicate title removed; Refs.~13, 14, 25 expanded to full author lists; Refs.~11, 19 volume/page numbers added; Ref.~20 CrCl$_3$ formula fixed; chemical formulae subscripted in Refs.~16, 18, 20--23, 31, 38--40, 60; Ref.~10 (arXiv) updated to journal version where available. The \texttt{achemso} BibTeX style enforces ACS formatting automatically.

\vspace{0.8em}

%--- R4C14 ---
\begin{reviewercomment}[Reviewer 4 --- Comment 14]
  14. Can the authors provide further justification for using $\Delta G_{\mathrm{ads}}$ as the primary descriptor for HER activity? Since HER is a multistep electrochemical process\ldots to what extent can $\Delta G_{\mathrm{ads}}$ alone be used to infer HER activity?
\end{reviewercomment}

\response We thank the Reviewer for this fundamental question. $\Delta G_{\mathrm{ads}}$ is the single most predictive thermodynamic descriptor for HER activity within the CHE framework (N{\o}rskov 2005) because: (1) the Volmer step (H adsorption) is rate-limiting in weak-binding regimes; (2) the Heyrovsky/Tafel step is limiting for strong binding; (3) both are captured by $\Delta G_{\mathrm{ads}}$ on the volcano plot. Full NEB kinetics is computationally prohibitive for 6 configurations and is identified as future work.

\manuscriptchange{\textbf{[Location: \textit{manuscript\_marked.pdf: Pages 29--30, Lines 542--569} \textbar\ \textit{manuscript\_clean.tex: Line 456} \textbar\ \textit{manuscript\_marked.tex: Line 475}]}\\ The discussion in Section~3.5 has been expanded (highlighted in blue) to explicitly acknowledge the thermodynamic scope and its connection to the volcano principle.}

\vspace{0.8em}

%--- R4C15 ---
\begin{reviewercomment}[Reviewer 4 --- Comment 15]
  15. Have the authors considered the thermodynamic stability of isolated Co, Fe, and Ni atoms on $\text{CrCl}_3$ against metal clustering or bulk-metal formation?\ldots
\end{reviewercomment}

\response We thank the Reviewer for this question. This is fully addressed in our response to Comment~3 above. The key result: Ni-embedded ($|E_{\mathrm{bind}}| = 4.95\text{ eV}$) is the only configuration thermodynamically stable vs. bulk Ni (4.44~eV/atom). All adsorbed configurations are kinetically stabilized single-site motifs.

\manuscriptchange{\textbf{[Location: \textit{manuscript\_marked.pdf: Pages 29--30, Lines 542--569} \textbar\ \textit{manuscript\_clean.tex: Line 456} \textbar\ \textit{manuscript\_marked.tex: Line 475}]}\\ See the manuscript addition described in our response to Comment~3 (Section~3.2, highlighted in blue).}

\vspace{0.8em}

%--- R4C16 ---
\begin{reviewercomment}[Reviewer 4 --- Comment 16]
  16. Why was the PBE functional without a Hubbard U correction selected\ldots have the authors assessed whether the relative ordering of Co, Fe, and Ni in terms of $\Delta G_{\mathrm{ads}}$ is robust with respect to the treatment of localized $d$ electrons?
\end{reviewercomment}

\response We thank the Reviewer for pressing on this point. This is addressed in our response to Comment~1. Additionally, the Co $<$ Fe $<$ Ni ordering of $\Delta G_{\mathrm{ads}}$ correlates with the H--surface interaction: Co ($-2.46\text{ eV}$) $>$ Fe ($-2.35\text{ eV}$) $>$ Ni ($-1.69\text{ eV}$). These quantities reflect the direct TM--H orbital overlap driven by TM $d$-electron count/valence, not by the $U$ value on Cr. DFT+U on TM dopants would be non-standard in this context and would require careful $U$ selection per metal.

\manuscriptchange{\textbf{[Location: \textit{manuscript\_marked.pdf: Page 7, Lines 145--147} \textbar\ \textit{manuscript\_clean.tex: Line 156} \textbar\ \textit{manuscript\_marked.tex: Line 175}]}\\ See the manuscript addition described in our response to Comment~1 (Section~2, highlighted in blue).}

\vspace{0.8em}

%--- R4C17 ---
\begin{reviewercomment}[Reviewer 4 --- Comment 17]
  17. Have the authors systematically investigated other possible adsorption sites for Co, Fe, and Ni besides the hollow site?\ldots
\end{reviewercomment}

\response We thank the Reviewer for this question. The hollow site is the canonical low-energy adsorption position for TM adatoms on CrCl$_3$ and related CrX$_3$ materials, consistent with Luo2020 (Li adsorption) and Mastrippolito2022 (O adsorption). The top-Cl site has the three Cl atoms too close for comfortable TM coordination; the top-Cr site places the TM directly above Cr with no Cl neighbors below. The hollow site provides the natural coordination pocket, and is the only position yielding stable relaxed configurations in our DFT optimizations.

\manuscriptchange{\textbf{[Location: \textit{manuscript\_marked.pdf: Page 16, Lines 331--335} \textbar\ \textit{manuscript\_clean.tex: Line 361} \textbar\ \textit{manuscript\_marked.tex: Line 380}]}\\ A statement confirming the hollow-site selection has been added to Section~3.2 (highlighted in blue in the marked manuscript).}

\vspace{0.8em}

%--- R4C18 ---
\begin{reviewercomment}[Reviewer 4 --- Comment 18]
  18. Can the authors distinguish more clearly between the thermodynamic and kinetic accessibility of the embedded and surface-adsorbed configurations?\ldots
\end{reviewercomment}

\response We thank the Reviewer for this question. The thermodynamic preference ($\Delta E_{\mathrm{emb-ads}}$) does not establish the kinetic barrier to incorporation. The AIMD trajectories provide partial kinetic information: Co and Fe show no incorporation within 5~ps at 300~K (suggesting a finite barrier), while Ni spontaneously migrates (indicating a low or absent barrier). NEB transition-state calculations for the migration pathway are identified as future work.

\manuscriptchange{\textbf{[Location: \textit{manuscript\_marked.pdf: Pages 20--21, Lines 402--412} \textbar\ \textit{manuscript\_clean.tex: Line 393} \textbar\ \textit{manuscript\_marked.tex: Line 412}]}\\ An explicit paragraph distinguishing thermodynamic preference from kinetic accessibility has been added to Section~3.2 (highlighted in blue): ``{\itshape \ExecuteMetaData[manuscript_marked.tex]{StabilityHierarchy}}''}

\vspace{0.8em}

%--- R4C19 ---
\begin{reviewercomment}[Reviewer 4 --- Comment 19]
  19. How relevant are the calculated adsorption properties of the proposed $\text{Co-CrCl}_3$ catalyst under realistic electrochemical conditions? Since the calculations are performed for an ideal, neutral monolayer in vacuum, how might solvation, electrolyte effects, applied potential, electric field, pH, surface charging, and possible structural reconstruction influence the predicted H adsorption thermodynamics?
\end{reviewercomment}

\response We thank the Reviewer for this important question. The CHE framework (N{\o}rskov 2005) maps the electrode potential onto the chemical potential of H$^+$/e$^-$, allowing vacuum DFT results to be referenced to the standard hydrogen electrode. Solvation of H$^*$ typically shifts $\Delta G_{\mathrm{ads}}$ by $-0.05$ to $-0.15\text{ eV}$ for $3d$ metal surfaces; applied potential enters via $-eU$; surface charging is secondary for covalently bound TM centers. These effects do not change the relative metal ordering. Implicit-solvent or constant-potential DFT methods are identified as future extensions.

\manuscriptchange{\textbf{[Location: \textit{manuscript\_marked.pdf: Page 32, Lines 610--612} \textbar\ \textit{manuscript\_clean.tex: Line 475} \textbar\ \textit{manuscript\_marked.tex: Line 494}]}\\ A statement has been added to the Conclusions (Section~4, highlighted in blue): ``\dots{\itshape \ExecuteMetaData[manuscript_marked.tex]{SolvationPotentialEffects}}''}

\vspace{0.8em}

%--- R4C20 ---
\begin{reviewercomment}[Reviewer 4 --- Comment 20]
  20. The statement that data will be made available on request is weak for a purely computational study. Please deposit the optimized structures (POSCAR or CIF files for all seven systems) in a public repository or provide them directly as Supporting Information.
\end{reviewercomment}

\response We thank the Reviewer for this recommendation. We have updated the Data Availability statement and included the optimized POSCAR files for all seven configurations in the Supporting Information.

\manuscriptchange{\textbf{[Location: \textit{manuscript\_marked.pdf: Page 33, Lines 634--639} \textbar\ \textit{manuscript\_clean.tex: Line 491} \textbar\ \textit{manuscript\_marked.tex: Line 510}]}\\ The Data Availability section has been updated (highlighted in blue): ``{\itshape \ExecuteMetaData[manuscript_marked.tex]{DataAvailabilityPOSCAR}}''}

\vspace{1.5em}
\noindent\rule{\textwidth}{0.6pt}

\newpage

%========================================================================================
%   RESPONSE TO REVIEWER 5
%========================================================================================
\section{Response to Reviewer \#5}

\noindent\textbf{Recommendation:} Major revisions needed as noted.

\vspace{0.5em}
\noindent\textbf{General Comments:}\\
\textit{The manuscript is well written, and the following key issues should be addressed before further consideration.}

\vspace{0.5em}
\noindent We thank Reviewer 5 for the constructive evaluation and recognition of the manuscript quality. We address all five points below.

\vspace{1em}

%--- R5C1 ---
\begin{reviewercomment}[Reviewer 5 --- Comment 1]
  1) The $\text{H}_2$ symbols in Fig. 2c are somewhat misleading and the subscript two means molecular hydrogen instead of the calculated H atom.
\end{reviewercomment}

\response We thank the Reviewer for this careful observation. The label has been corrected from $\text{H}_2$ to atomic H (and H$^*$ where adsorbed) in Figure~2(c) and throughout the manuscript. The figure caption for Figure~2 explicitly clarifies: ``$\Delta G_{\mathrm{ads}}=1.82$~eV'' refers to a single adsorbed H atom referenced to $\frac{1}{2}\text{H}_2$, consistent with the CHE convention.

\manuscriptchange{\textbf{[Location: \textit{manuscript\_marked.pdf: Page 15, Lines 308--314} \textbar\ \textit{manuscript\_clean.tex: Line 351} \textbar\ \textit{manuscript\_marked.tex: Line 370}]}\\ Figure~2(c) labels corrected to H$^*$ / H$_{\mathrm{ads}}$ throughout. The Figure~2 caption has been updated (highlighted in blue).}

\vspace{0.8em}

%--- R5C2 ---
\begin{reviewercomment}[Reviewer 5 --- Comment 2]
  2) Can the authors reproduce the reported hydrogen adsorption free energy of $0.13\text{ eV}$ in Ref. 40 by using the same models and calculation parameters as the literature? For the two rather different values, $0.13$ (Ref. 40) and $1.82\text{ eV}$ (this work), which one represent the real HER performance of pristine monolayer $\text{CrCl}_3$? More discussions should be added here to clarify this important point.
\end{reviewercomment}

\response We thank the Reviewer for emphasizing this important clarification. As detailed in our response to Reviewer~3 (Comment~1), the two values correspond to \textbf{geometrically and chemically distinct} adsorption regimes:
\begin{itemize}[leftmargin=1.5em, itemsep=0.2em]
  \item Our $1.82\text{ eV}$: H chemisorbed directly at the Cr center (site S$_3$, $d_{\mathrm{H-Cr}} = 1.55\text{ \AA}$), the only locally stable chemisorbed state under our protocol.
  \item Zala et al. $0.13\text{ eV}$: H in a weakly interacting on-Cl or near-Cl surface configuration with different coverage, supercell, and computational parameters.
\end{itemize}
Both confirm that pristine CrCl$_3$ has poor-to-mediocre HER activity. Our computed total energies from the VASP OUTCAR files verify the DFT energetics:
\begin{itemize}[leftmargin=1.5em, itemsep=0.2em]
  \item $E(\text{H}_2) = -6.760\text{ eV}$ ($\frac{1}{2}E_{\text{H}_2} = -3.380\text{ eV}$), calculated with the same ENCUT = 400 eV and EDIFF = $10^{-6}\text{ eV}$ parameters used throughout.
  \item The electronic contribution $\Delta E_{\mathrm{H}}$ for each doped system is computed consistently from the same reference H$_2$ energy.
\end{itemize}

\manuscriptchange{\textbf{[Location: \textit{manuscript\_marked.pdf: Pages 15--16, Lines 315--324} \textbar\ \textit{manuscript\_clean.tex: Line 355} \textbar\ \textit{manuscript\_marked.tex: Line 374}]}\\ An expanded comparison has been added to Section~3.1 (highlighted in blue), as quoted in our response to Reviewer~3, Comment~1.}

\vspace{0.8em}

%--- R5C3 ---
\begin{reviewercomment}[Reviewer 5 --- Comment 3]
  3) Why only the hollow sites are considered for the transition metal doping? Have they tested the doping at other sites (e.g., the S1 and S3 sites in Fig. 2)? How about their energetic stability compared to the hollow site?
\end{reviewercomment}

\response We thank the Reviewer for this question. We systematically tested initial transition-metal placements at top-Cl (S$_1$), hollow (S$_2$), and top-Cr (S$_3$) positions:
\begin{itemize}[leftmargin=1.5em, itemsep=0.2em]
  \item TM atoms initialized at top-Cl (S$_1$) or top-Cr (S$_3$) spontaneously relaxed into the threefold hollow site (S$_2$), which maximizes coordination to three surface Cl anions ($d_{\mathrm{TM-Cl}} \approx 2.16\text{--}2.26\text{ \AA}$). This was verified from our fully relaxed CONTCAR files.
  \item The hollow site provides the natural coordination pocket defined by three Cl atoms, and is the only position yielding stable relaxed configurations for Co, Fe, and Ni.
  \item This is consistent with the established literature: Luo2020 (Li at hollow on CrCl$_3$), Mastrippolito2022 (O at hollow).
\end{itemize}

\manuscriptchange{\textbf{[Location: \textit{manuscript\_marked.pdf: Page 16, Lines 331--335} \textbar\ \textit{manuscript\_clean.tex: Line 361} \textbar\ \textit{manuscript\_marked.tex: Line 380}]}\\ A statement confirming the hollow-site selection and other-site testing has been added to Section~3.2 (highlighted in blue).}

\vspace{0.8em}

%--- R5C4 ---
\begin{reviewercomment}[Reviewer 5 --- Comment 4]
  4) Besides the binding energies, dissolve potentials should be further evaluated to confirm the metal doping stability under the electrochemical condition.
\end{reviewercomment}

\response We thank the Reviewer for this important electrochemical consideration. The strong covalent TM--Cl binding energies ($-2.60$ to $-3.38\text{ eV}$ for adsorbed; $-3.36$ to $-4.95\text{ eV}$ for embedded) indicate robust local anchoring. Under acidic aqueous conditions ($\text{M} \rightarrow \text{M}^{2+} + 2e^-$), these binding energies shift the effective dissolution potential positively relative to bulk metal dissolution. Combined with the ICOHP-confirmed covalent bonding (cumulative TM--Cl ICOHP of $-6.18$ to $-6.84\text{ eV}$ for adsorbed; $-10.01$ to $-11.40\text{ eV}$ for embedded), the TM centers are covalently anchored rather than weakly adsorbed. A formal Pourbaix stability analysis is identified as a direction for future electrochemical modeling.

\manuscriptchange{\textbf{[Location: \textit{manuscript\_marked.pdf: Page 17, Lines 353--359} \textbar\ \textit{manuscript\_clean.tex: Line 372} \textbar\ \textit{manuscript\_marked.tex: Line 391}]}\\ A discussion of dissolution stability and Pourbaix considerations has been added to Section~3.2 (highlighted in blue).}

\vspace{0.8em}

%--- R5C5 ---
\begin{reviewercomment}[Reviewer 5 --- Comment 5]
  5) For each of Co, Fe, and Ni, which doping configuration (surface vs. embedding) in Fig. 3 is energetically favorable? The better one should be used for HER if the energy difference is substantial.
\end{reviewercomment}

\response We thank the Reviewer for this important question. The relative thermodynamic stability at 0~K is:
\begin{center}
  \small
  \begin{tabular}{lcl}
    \toprule
    Metal & $\Delta E_{\mathrm{emb-ads}}$ (eV) & Thermodynamically favored at 0~K \\
    \midrule
    Co    & $-0.73$                            & Embedded                         \\
    Fe    & $-0.76$                            & Embedded                         \\
    Ni    & $-1.57$                            & Embedded                         \\
    \bottomrule
  \end{tabular}
\end{center}

However, from an electrocatalytic standpoint, the embedded motif is catastrophic for HER: $\Delta G_{\mathrm{ads}} = 1.75\text{ eV}$ (Co-embedded), $1.36\text{ eV}$ (Fe-embedded), and $3.62\text{ eV}$ (Ni-embedded) -- all far from thermoneutral. In contrast, surface-adsorbed Co delivers $\Delta G_{\mathrm{ads}} = 0.20\text{ eV}$ ($\eta = 0.20\text{ V}$). At room temperature (300~K), AIMD confirms that surface-adsorbed Co and Fe remain kinetically trapped in the surface state (Figure~3(g,h)), preserving the catalytically active, open sites. Therefore, while all three metals thermodynamically prefer embedding (especially Ni, which spontaneously incorporates during the 300~K AIMD trajectory), the kinetically accessible and catalytically optimal motif for near-thermoneutral HER is the surface-adsorbed configuration of Co.

\manuscriptchange{\textbf{[Location: \textit{manuscript\_marked.pdf: Pages 20--21, Lines 402--412} \textbar\ \textit{manuscript\_clean.tex: Line 393} \textbar\ \textit{manuscript\_marked.tex: Line 412}]}\\ A discussion highlighting this fundamental stability--activity trade-off has been added to Section~3.4 (highlighted in blue in the marked manuscript). The Conclusions also state: ``\textit{These results identify surface-adsorbed Co as the most favorable H-binding motif among the configurations examined and establish coordination engineering as a promising strategy for tuning the reactivity of magnetic two-dimensional materials.}''}

\end{document}